\documentclass[10pt,a4paper]{amsart}
\usepackage[margin=19mm]{geometry}
\usepackage{amsmath,amssymb,mathtools}
\usepackage[scr=rsfs,cal=euler]{mathalfa}
\usepackage{xfrac}
\usepackage[unicode,hidelinks,bookmarks=false]{hyperref}
\newcommand{\ie}{i.e.\ }
\newcommand{\cf}{cf.\ }
\newcommand{\resp}{respectively\ }
\newcommand{\Subsection}{\subsection}
\providecommand{\nobreakdash}{\nobreak\hbox{-}}
\setlength{\emergencystretch}{2em}
\multlinegap0pt
\usepackage{bm}
\usepackage{bbm}
\usepackage{dsfont}
\usepackage{xspace}
\usepackage{cancel}
\usepackage{mathtools}
\usepackage{amsthm}
\makeatletter
\@ifundefined{claim}{\newtheorem{claim}{Claim}}{}
\@ifundefined{lemma}{\newtheorem{lemma}{Lemma}}{}
\makeatother
\newcommand{\ex}{\mathrm{ex}}
\title[Planar Tur\'an number of disjoint union of $C\_3$ and $C\_5$]{Planar Tur\'an number of disjoint union of $C\_3$ and $C\_5$}
\author{Luyi Li, Ping Li, Guiying Yan, Qiang Zhou}
\date{Source: arXiv:2507.16351v1 (CC BY 4.0)}
\begin{document}
\maketitle

\noindent\emph{Source and licence.} Planar Tur\'an number of disjoint union of $C_3$ and $C_5$. Version arXiv:2507.16351v1 (CC BY 4.0), distributed under the Creative Commons Attribution 4.0 International licence, as recorded in the arXiv licence metadata for this exact version and shown on the abstract page https://arxiv.org/abs/2507.16351v1. Full rights notice: licenses/arxiv-2507.16351-CC-BY-4.0.txt.\par\smallskip

\noindent\textit{Editorial scope.} Counted unit: the Lemma whose source label is \texttt{lemma1} in the article (source0.tex lines 108--119, source label \texttt{lemma1}), delivered together with the article's own complete original proof of it (source0.tex lines 121--146, one \texttt{proof} environment of 26 lines). No mathematics is written, repaired or re-derived: statement and proof are the article's own text line for line. Required context retained uncounted: none. Every definition, display and label of those lines is retained unchanged. Omitted from this package and not redistributed: 1--107, 120--120, 147--339. Nothing inside a retained excerpt is omitted or altered: every retained body is delivered exactly as the article's lines are written, so no omission receipt of any kind accompanies this package. Citations: the retained excerpts carry no citation command at all, so no bibliography entry of the article is transcribed and no reference is redistributed. Reference adaptation: none was needed and none was made; every reference inside the delivered unit resolves inside it, so no unresolved reference or citation is produced. Printed slips kept literal: no printed word, symbol, hypothesis or number of the counted statement or of its proof was altered. Layout and class: the article's own class is \texttt{article} with packages including geometry, float, latexsym, bm, amsmath, amsfonts, amsthm, amssymb, bbding, mathrsfs, mathtools, fancyhdr, graphicx, hyperref; the wrapper is the lane's pinned \texttt{amsart} editorial wrapper at 10pt on A4, which restates the theorem-like environments the retained text needs and adds the article's own macro definitions that the retained text needs (\texttt{\textbackslash ex}), with extra wrapper packages (bm, bbm, dsfont, xspace, cancel, mathtools). The delivered numbering is the unit's own. Assets: no retained span contains a figure environment, a graphics-inclusion command, a PDF-page inclusion, a table or a TikZ picture; no image file is copied into this package, the package declares no asset dependency and no image file is redistributed with it. Licence: version arXiv:2507.16351v1 of this article is distributed under the Creative Commons Attribution 4.0 International licence, as recorded in the arXiv licence metadata for this exact version and shown on the abstract page https://arxiv.org/abs/2507.16351v1. A full rights notice is delivered at licenses/arxiv-2507.16351-CC-BY-4.0.txt. Characters: every retained excerpt is pure ASCII, so the delivered engine, XeLaTeX with its default Latin Modern text font, reports no missing character.
\begin{lemma}\label{lemma1}
    Let $G$ be a plane graph with $e(G)\geq \left\lfloor\frac{8n-13}{3}\right\rfloor$. If $G$ is $C_3\cup C_5$-free and $n\geq 295660$, then one of the following statements holds:
    \begin{enumerate}
        
        \item[(1)] $G$ has at least $3$ triangular blocks containing $C_5$ and intersecting exact two vertices;
        
        \item[(2)] $G$ has a triangular block with at least $521$ vertices;
        
        \item[(3)] $G$ has at least $4$ triangular blocks containing $C_5$ and intersecting exact one vertex.
    
    \end{enumerate}
\end{lemma}

\begin{proof}
    Assume that statements (2) and (3) do not hold. Next we prove that statement (1) holds. Without loss of generality, suppose that $G$ is the planar embedding.
    Since $e(G)>\ex_{\mathcal{P}}(n,C_5)=\frac{12n-33}{5}$ when $n\geq 11$, $G$ contains a copy $D$ of $C_5$. Now we have the following claims:
    
    \begin{claim}\label{claim1}
        There are at most $10$ triangular blocks isomorphic to $B_4^2$ in $G$.
    \end{claim}
    \begin{proof}
        Note that each $B_4^2$ intersects at least two vertices of $D$. Otherwise, we could easily find $C_3\cup C_5$. Since any two blocks of $B_4^2$ intersect at most one vertex, it follows that there are at most $10$ triangular blocks isomorphic to $B_4^2$ in $G$.
    \end{proof}
    
    \begin{claim}\label{claim2}
        There are at least $f(n)$ triangular blocks containing $C_5$, where $f(n)=\frac{4n+15097}{15555}$.
    \end{claim}
    \begin{proof}
        Assume that there are only $\alpha=f(n)-1$ triangular blocks containing $C_5$. By Claim~\ref{claim1}, there are at most $10$ triangular blocks isomorphic to $B_4^2$ in $G$. By the assumption that statement (2) does not hold, the largest triangular block has at most $520$ vertices. The remaining triangular blocks that do not contain $C_5$ and are not isomorphic to $B_4^2$, denoted as $B_1,\ldots,B_t$, are isomorphic to $B_2^1$, $B_3^1$, and $B_4^1$. Then
        $$f_3(G) \leq \sum_{i\in[t]}f_3(B_i)+3\times10+\alpha\times (2\times 520-3) \leq \frac{2}{5}\sum_{i\in [t]}e(B_i)+1037\alpha+30 \leq \frac{2}{5}e(G)+1037\alpha+30,$$
        By Euler's formula, we have
        $$2e = \sum_{i}if_i \geq 3f_3+4(f-f_3) \geq 4(e+2-n)-\frac{2e}{5}-1037\alpha-30,$$
        which implies $e \leq \frac{5}{8}(4n+1037\alpha+22)=\frac{8n-16}{3}$, thus $e<\left\lfloor\frac{8n-13}{3}\right\rfloor$, a contradiction.   
    \end{proof}
    Since $G$ is $C_3\cup C_5$-free, each triangular block containing $C_5$ intersects with $D$.
    Recall that statement (3) does not hold. It follows that there are at most $3$ triangular blocks containing $C_5$ and intersecting exact one fixed vertex with $D$. Since $G$ is a plane graph, for any subset $A\subseteq V(D)$ with $|A|\geq 3$, there are at most two triangular blocks $B_1,B_2$ such that $B_i$ contains a $5$-cycle and $B_i\cap V(D)=A$ for $i=1,2$. Thus, there are at least 
    $$\frac{f(n)-3\times 5-2\big[\binom{5}{3}+\binom{5}{4}+\binom{5}{5}\big]}{\binom{5}{2}}=\frac{2n-357994}{77775}\geq 3$$ 
    triangular blocks that contain $C_5$ and intersect exactly two vertices. The result follows.
\end{proof}

\end{document}
